\subsection{Homological characterization of regular Noetherian rings}

THEOREM 1 (Serre).--- For a local Noetherian ring to be regular, it is necessary and sufficient that its homological dimension be finite.

We have seen that a regular local Noetherian ring has finite homological dimension (prop. 1).

Conversely, let A be a local Noetherian ring of finite homological dimension n; by § 3, n° 3, cor. 2 of prop. 4 and n° 5, th. 1, we have

\[
n = \mathrm{dh} (\mathrm{A}) = \mathrm{dp} _ {\mathrm{A}} (\kappa_ {\mathrm{A}}) = \operatorname{prof} (\mathrm{A}).
\]

If \(n=0\) , the A-module \(\kappa_{A}\) is free, hence \(\mathfrak{m}_{A}=0\) and A is a field. Suppose \(n>0\) and reason by induction on \(n\) . Since \(\operatorname{prof}(\Lambda)>0\) , the ideal \(\mathfrak{m}_{A}\) is not associated with A ( \(\S 1\) , n° 1, remark 2), hence is not contained in the union of \(\mathfrak{m}_{A}^{2}\) and the ideals associated with A (II, \(\S 1\) , n° 1, prop. 2). Consequently (IV, \(\S 1\) , n° 1, cor. 2 of prop. 2), one can find an element \(x\) of \(\mathfrak{m}_{A}-\mathfrak{m}_{A}^{2}\) such that the homothety \(x_{A}\) is injective. Let B denote the local Noetherian ring A/ \(xA\) and consider the sequence of A-modules

\[
0 \to \kappa_ {\mathrm{A}} \stackrel {i} {\longrightarrow} \mathfrak {m} _ {\Lambda} / x \mathfrak {m} _ {\Lambda} \stackrel {p} {\longrightarrow} \mathfrak {m} _ {\mathrm{B}} \to 0
\]

where the map \(i\) is deduced by passage to quotients from the map \(a \mapsto ax\) of A into \(\mathfrak{m}_{\mathrm{A}}\) and where \(p\) is the canonical surjection; it is exact. Since the class of \(x\) in the \(\kappa_{\mathrm{A}}\) -vector space \(\mathfrak{m}_{\mathrm{A}} / \mathfrak{m}_{\mathrm{A}}^{2}\) is nonzero, there exists an A-linear map \(\phi : \mathfrak{m}_{\mathrm{A}} \to \kappa_{\mathrm{A}}\) with \(\phi(x) = 1\) ; by passage to quotients, one deduces from \(\phi\) a retraction of \(i\) , so that the preceding exact sequence is split. This entails the relations

\[
\mathrm{dp} _ {\mathrm{B}} \left(\mathfrak {m} _ {\mathrm{B}}\right) \leqslant \mathrm{dp} _ {\mathrm{B}} \left(\mathfrak {m} _ {\mathrm{A}} / x \mathfrak {m} _ {\mathrm{A}}\right) = \mathrm{dp} _ {\mathrm{A}} \left(\mathfrak {m} _ {\mathrm{A}}\right) <   + \infty
\]

(cor. 2 of prop. 7 of § 3, n° 4 and A, X, p.~135, cor. 1). Cor. 2 of loc. cit. applied to the exact sequence of B-modules 0 → mB → B → κB → 0 entails dpB(κB) \textless{} +∞. The ring B is therefore of finite homological dimension (§ 3, n° 3, cor. 2 of prop. 4), and of depth n - 1 (§ 1, n° 4, prop. 7 and n° 3, cor. of prop. 4). It follows from the induction hypothesis that B is regular, hence that A is regular (VIII, § 5, n° 3, cor. 1 of prop. 2).

Consequently, if A is a local Noetherian ring, there is an equivalence between the following three properties:

\begin{enumerate}
\def\labelenumi{(\roman{enumi})}
\item
  A is regular;
\item
  the A-module \(\kappa_{A}\) has finite projective dimension;
\item
  every finitely generated A-module has projective dimension \(<+\infty\) .
\end{enumerate}

DEFINITION 1.--- A ring A is said to be regular if it is Noetherian and the local ring \(\mathrm{A}_{\mathfrak{m}}\) is regular for every maximal ideal \(\mathfrak{m}\) of \(\mathrm{A}\) .

PROPOSITION A.--- Let A be a Noetherian ring. The following conditions are equivalent:

\begin{enumerate}
\def\labelenumi{(\roman{enumi})}
\item
  A is regular;
\item
  every finitely generated A-module has projective dimension \(< +\infty\) ;
\item
  for every maximal ideal m of A, the projective dimension of A/m is finite;
\item
  for every prime ideal p of A, the local ring A \(_{p}\) is regular.
\end{enumerate}

Let p be a prime ideal of A; if the A-module A/p has finite projective dimension, then so does the \(A_{p}\) -module \(\kappa(\mathfrak{p})\) by cor. 1 of prop. 3 of § 3, n° 2, so that the local ring \(A_{p}\) is regular (th. 1). It follows that (ii) implies (iv) and that (iii) implies (i). The implications (iv)⇒(i) and (ii)⇒(iii) are clear.

Let us prove that (i) implies (ii). Let M be a finitely generated A-module. Under hypothesis (i), we have \(\mathrm{dp}_{\mathrm{A}_{\mathfrak{m}}}(M_{\mathfrak{m}}) \leqslant \mathrm{dh}(A_{\mathfrak{m}}) < +\infty\) for every maximal ideal \(\mathfrak{m}\) of A ( \(n^{\circ} 1\) , prop. 1); hence M has finite projective dimension \(< +\infty\) ( \(\S 3, n^{\circ} 2\) , cor. 2 of prop. 3), whence (ii).

Examples.--- 1) If the ring A is regular, the ring of fractions S \(^{-1}\) A is regular for every multiplicative subset S of A: this follows for example from characterization (iii) above.

\begin{enumerate}
\def\labelenumi{\arabic{enumi})}
\setcounter{enumi}{1}
\item
  For a ring to be regular it is necessary and sufficient that it be isomorphic to the product of a finite family of regular integral domains; this follows in fact from the fact that every regular ring is locally integral (§ 1, no. 8), since regular local rings are integral domains.
\item
  The regular integral domains of dimension \(\leqslant 1\) are the Dedekind rings (VIII, § 5, no. 1, example 1 and VII, § 2, no. 2, theorem 1).
\end{enumerate}

COROLLARY 1. Let A be a Noetherian ring. The following conditions are equivalent:

\begin{enumerate}
\def\labelenumi{(\roman{enumi})}
\item
  ; we have \(\mathrm{dh(A)} < +\infty\)
\item
  A is regular and one has \(\dim(A) < +\infty\) .
\end{enumerate}

If these conditions are satisfied, one has \(\dim(\Lambda) = \mathrm{d}\mathrm{h}(\mathrm{A})\) .

If the ring A is regular, one has for every maximal ideal m of Λ the equality \(\dim(A_{\mathfrak{m}}) = \mathrm{dh}(A_{\mathfrak{m}})\) (no. 1, prop. 1), and hence

\[
\mathrm{dh} (\mathrm{A}) = \sup _ {\mathfrak {m}} \mathrm{dh} (\mathrm{A} _ {\mathfrak {m}}) = \sup _ {\mathfrak {m}} \dim (\mathrm{A} _ {\mathfrak {m}}) = \dim \mathrm{A}
\]

(§ 3, no. 2, prop. 3 and VIII, § 1, no. 3, prop. 8). On the other hand, if dh(Λ) \textless{} +∞, the ring A is regular by prop. 4. The corollary follows.

There exist regular Noetherian rings of infinite dimension (VIII, § 5, exerc. 6).

COROLLARY 2. A regular ring is normal, Gorenstein and Macaulay.

Indeed, a regular local ring is integrally closed (VIII, § 5, no. 2, cor. 1 of th. 1), Gorenstein (§ 3, no. 9, example 4) and Macaulay (§ 2, no. 5, example 6).

COROLLARY 3.--- Let A be a Noetherian ring, J an ideal of A and  the separated completion of A for the J-adic topology.

\begin{enumerate}
\def\labelenumi{\alph{enumi})}
\item
  For the ring \(\widehat{\mathbf{A}}\) to be regular, it is necessary and sufficient that, for every maximal ideal \(\mathfrak{m}\) of \(\Lambda\) containing J, the ring \(\mathrm{A}_{\mathfrak{m}}\) be regular.
\item
  If the ring A is regular, the ring \(\hat{\Lambda}\) is regular. If the ring \(\hat{\mathrm{A}}\) is regular and the ideal J is contained in the radical of A, the ring \(\Lambda\) is regular.
\end{enumerate}

By prop. 8 of III, § 3, no. 4, for the ring \(\widehat{\Lambda}\) to be regular, it is necessary and sufficient that it be so for \(\widehat{A}_{\widehat{\mathfrak{m}}}\) for every maximal ideal \(\mathfrak{m}\) of A containing J. Since the completions of the local rings \(\widehat{A}_{\widehat{\mathfrak{m}}}\) and \(A_{\mathfrak{m}}\) are isomorphic (loc. cit.), assertion a) follows from VIII, § 5, no. 1, cor. of prop. 1. Assertion b) follows from a).

COROLLARY 4.--- Let A be a regular ring and P a finitely generated projective A-module. The symmetric algebra \(\mathsf{S}_{\mathrm{A}}(\mathrm{P})\) is a regular ring.

Let p be a prime ideal of \(\mathbf{S}_{\mathrm{A}}(\mathrm{P})\) and q its inverse image in A. The local ring \(\mathbf{S}_{\mathrm{A}}(\mathrm{P})_{\mathfrak{p}}\) is a ring of fractions of the ring \(\mathbf{S}_{\Lambda}(\mathrm{P})_{\mathfrak{q}}\) , which is isomorphic to \(\mathbf{S}_{\mathrm{A}_{q}}(\mathrm{P}_{\mathfrak{q}})\) (A, III, p.~72, prop. 7); it suffices to prove that the latter is regular. This reduces us to the case where A is local; but then P is free of finite type. By prop. 1 of no. 1 and A, X, p.~143, cor. 1, one has \(\mathrm{dh}(\mathbf{S}_{\mathrm{A}}(\mathrm{P})) = \mathrm{dh}(\mathrm{A}) + \mathrm{rg}_{\mathrm{A}}(\mathrm{P}) < +\infty\) , and by \(\mathbf{S}_{\mathrm{A}}(\mathrm{P})\) is regular by prop. 4.

COROLLARY 5. Let A be a regular ring, and \((\mathrm{T}_i)_{i\in I}\) a finite family of indeterminates. The polynomial ring \(\mathrm{A}[(\mathrm{T}_i)_{i\in I}]\) and the formal power series ring \(\mathrm{A}[(\mathrm{T}_i)_{i\in I}]\) are regular.

This follows from cor. 4 and cor. 3 b).

